% =============================================================
% Список обозначений и сокращений
% =============================================================
\chapter*{Список обозначений и сокращений}
\phantomsection\addcontentsline{toc}{chapter}{Список обозначений и сокращений}

\section*{Латинские обозначения}

\begin{tabularx}{\textwidth}{@{}p{3.5cm}X@{}}
$a$, $b$ & полуоси элемента микрорельефа, м \\
$c$ & радиальный зазор подшипника, м \\
$C_{ij}$ & коэффициенты демпфирования смазочного слоя, Н$\cdot$с/м \\
$D$ & диаметр подшипника, м \\
$e$ & эксцентриситет вала, м \\
$f$ & коэффициент трения \\
$F$ & несущая сила (нагрузочная способность), Н \\
$F_f$ & сила трения, Н \\
$h$ & толщина смазочного слоя (функция зазора), м \\
$h_{\min}$ & минимальная толщина смазочного слоя, м \\
$h_p$ & глубина элемента микрорельефа, м \\
$H$ & безразмерная толщина смазочного слоя \\
$H_p$ & безразмерная глубина элемента микрорельефа \\
$K$ & параметр клина упорного подшипника \\
$K_{ij}$ & коэффициенты жёсткости смазочного слоя, Н/м \\
$L$ & длина подшипника, м \\
$L_1$, $L_2$ & характерные масштабы длины, м \\
$M_f$ & момент трения, Н$\cdot$м \\
$n$ & частота вращения, об/мин \\
$N$ & число колодок упорного подшипника \\
$p$ & давление в смазочном слое, Па \\
$p_0$ & характерный масштаб давления, Па \\
$p_{\mathrm{cav}}$ & давление кавитации, Па \\
$P$ & безразмерное давление \\
$Q$ & расход смазочного материала, м$^3$/с \\
$r$ & радиальная координата, м \\
$R$ & радиус вала (цапфы), м \\
$R_a$ & среднеарифметическое отклонение профиля шероховатости, м \\
$R_q$ & среднеквадратическая шероховатость, м \\
$S$ & число Зоммерфельда \\
$S_p$ & относительная площадь текстурирования \\
$T$ & температура, К \\
$U$ & скорость скольжения, м/с \\
$W$ & несущая сила упорного подшипника, Н \\
$x$, $y$, $z$ & декартовы координаты, м \\
$Z$ & безразмерная осевая координата \\
\end{tabularx}

\section*{Греческие обозначения}

\begin{tabularx}{\textwidth}{@{}p{3.5cm}X@{}}
$\alpha$ & угол ориентации элемента микрорельефа, рад \\
$\beta$ & угол сектора колодки упорного подшипника, рад \\
$\varepsilon$ & относительный эксцентриситет ($\varepsilon = e/c$) \\
$\vartheta$ & степень заполнения (в модели кавитации JFO) \\
$\theta$ & окружная координата, рад \\
$\lambda_r$ & параметр смазочного слоя ($\lambda_r = h_{\min}/R_q$) \\
$\Lambda$ & параметр анизотропии уравнения Рейнольдса ($\Lambda = L_1/L_2$) \\
$\mu$ & динамическая вязкость смазки, Па$\cdot$с \\
$\rho$ & плотность смазки, кг/м$^3$ \\
$\phi$ & коэффициент покрытия (доля текстурированной площади) \\
$\phi_p$, $\phi_s$ & коэффициенты потоков Патира--Ченга \\
$\varphi$ & окружная координата радиального подшипника, рад \\
$\psi$ & относительный зазор ($\psi = c/R$) \\
$\omega$ & угловая скорость вращения, рад/с \\
$\omega_{\mathrm{SOR}}$ & параметр релаксации метода SOR \\
\end{tabularx}

\section*{Сокращения}

\begin{tabularx}{\textwidth}{@{}p{3.5cm}X@{}}
GPU & графический процессор (\textit{Graphics Processing Unit}) \\
JFO & модель кавитации Якобссона--Флоберга--Олссона \\
SOR & метод последовательной верхней релаксации (\textit{Successive Over-Relaxation}) \\
THD & термогидродинамический (\textit{Thermohydrodynamic}) \\
ДМ & детерминированные микроструктуры \\
СФУ & Сибирский федеральный университет \\
СТУ & стандарт университета \\
\end{tabularx}

\section*{Индексы}

\begin{tabularx}{\textwidth}{@{}p{3.5cm}X@{}}
cav & кавитационный \\
ext & внешний \\
in & входной \\
max & максимальный \\
min & минимальный \\
out & выходной \\
smooth & гладкая поверхность \\
tex & текстурированная поверхность \\
\end{tabularx}

\newpage
